\section{Hoare Logic for an Untyped $\lambda$-Calculus with Recursion}
\label{sec:lam}

For IMP, step-indexing was overkill: the index was a presentation
convenience.  We now move to a language where it is no longer
optional --- an untyped $\lambda$-calculus with a fix-point operator.
Two things change.  The wp predicate must explicitly encode
\emph{safety}, since the language has stuck terms (e.g.\
$\mathtt{EApp}\, (\mathtt{EInt}\, 5)\, (\mathtt{EInt}\, 3)$); and the
recursion introduced by \texttt{fix} turns the soundness predicate
for a recursive function into something defined in terms of itself,
so that the natural soundness proof no longer goes through meta-level
induction on the step index --- it goes through the internal Löb
rule, used directly, with no outer induction on $n$.

\paragraph{The language.}
A de-Bruijn representation with integer constants, an \texttt{if0}
conditional, $\lambda$-abstraction and application, and an explicit
fix-point binder:
\begin{lstlisting}
Inductive expr : Type :=
  | EVar  : nat -> expr  | EInt  : Z -> expr
  | EPlus : expr -> expr -> expr
  | EIf0  : expr -> expr -> expr -> expr
  | ELam  : expr -> expr | EApp  : expr -> expr -> expr
  | EFix  : expr -> expr.
\end{lstlisting}
\noindent
$\mathtt{EFix}\, body$ binds the self-reference at de-Bruijn index
$0$; we treat it as a value so it can be passed around, with
self-unfolding firing only on application.  The small-step relation
includes the usual arithmetic and conditional rules, $\beta$-reduction
$\mathtt{EApp}\, (\mathtt{ELam}\, body)\, v \to body[v / 0]$, and the
fix-unfolding rule
\begin{lstlisting}
| StepFix : forall body v, value v ->
    EApp (EFix body) v ~> EApp (subst (EFix body) body) v.
\end{lstlisting}
\noindent
which substitutes the self-reference before letting an ordinary
$\beta$-reduction proceed.  Reductions are call-by-value.

\paragraph{Step-indexed wp with safety.}
The wp predicate adds an explicit safety conjunct:
\begin{lstlisting}
Fixpoint wp_at (n : nat) (e : expr) (Q : expr -> Prop) : Prop :=
  match n with
  | O => True
  | S k => (value e -> Q e) /\
           (value e \/ exists e', step e e') /\
           (forall e', step e e' -> wp_at k e' Q)
  end.
\end{lstlisting}
\noindent
The middle conjunct rules out stuck non-value expressions, which
would otherwise satisfy the wp vacuously (the first conjunct is
discharged by non-valuehood, and the third quantifies over an empty
step relation).  Two general rules carry every soundness proof:
\texttt{wp\_value}, saying a value satisfies the wp targeting any $Q$
it semantically satisfies, and the pure-step rule \texttt{wp\_pure\_step\_at},
saying that if every reduct of $e$ is the same $e'$ and $e'$
satisfies the wp at index $n$, then $e$ satisfies it at $n+1$.  The
pure-step rule propagates the wp through any $\beta$-reduction or
fix-unfolding.  We do not develop dedicated rules for \texttt{EPlus},
\texttt{EIf0}, or \texttt{EApp}; they either need an evaluation-context
framework or unfold case-by-case via \texttt{wp\_pure\_step}, and
neither is necessary for the section's conceptual point.

\paragraph{The specification predicate for recursive functions.}
The simplest non-trivial spec shape is the \emph{closed} one, where a
single predicate $S$ plays both pre- and postcondition:
\begin{lstlisting}
Definition recursive_spec (e : expr) (S : expr -> Prop) : iProp :=
  MkProp (fun n => forall v, value v -> S v -> wp_at n (EApp e v) S)
         (recursive_spec_mono e S).
\end{lstlisting}
\noindent
$\mathtt{recursive\_spec}\, e\, S$ at depth $n$ says: ``for every
value $v$ satisfying $S$, applying $e$ to $v$ is safe through $n$
steps and produces a value satisfying $S$.''  Crucially, $S$ appears
twice; this is the source of the self-reference.  To prove the spec
for $\mathtt{EFix}\, body$ at depth $n+1$ we need to know it at depth
$n$, in order to handle the recursive call inside the body.

\paragraph{The \texttt{wp\_fix} soundness rule.}
\begin{theorem}[\texttt{wp\_fix}]
For any \texttt{body} and predicate $S$, suppose that for every value
$e_{\mathrm{self}}$ satisfying $\mathtt{recursive\_spec}\, e_{\mathrm{self}}\, S$,
we have $\mathtt{recursive\_spec}\, (\mathtt{subst}\, e_{\mathrm{self}}\, body)\, S$.
Then $\top \vdash \mathtt{recursive\_spec}\, (\mathtt{EFix}\, body)\, S$.
\end{theorem}
\noindent
That is, if the body is a well-behaved transformer of recursive specs
--- producing a function satisfying the spec from a self-reference
satisfying the spec --- the fix-point itself satisfies the spec.  The
natural proof factors the goal through Löb's form via
\texttt{ientails\_trans} and closes with \texttt{apply iLob}; the
finite intermediate $\top \vdash \triangleright \varphi \to \varphi$
case-analyses on the depth supplied to the \texttt{iImpl}, uses
\texttt{wp\_pure\_step\_at} for the fix-unfolding step, and discharges
the recursive obligation from the $\mathtt{Hlater}$ supplied at the
predecessor index:
\begin{lstlisting}
Theorem wp_fix body S :
  (forall e_self, value e_self ->
     recursive_spec e_self S
       ⊢ recursive_spec (subst e_self body) S) ->
  ⊤ ⊢ recursive_spec (EFix body) S.
Proof.
  intros Hbody.
  apply ientails_trans with
    (iLater (recursive_spec (EFix body) S)
       → recursive_spec (EFix body) S).
  - intros n _ m Hmn Hlater.
    destruct m as [|m'].
    + intros v Hvalue HS. simpl. trivial.
    + simpl in Hlater. intros v Hvalue HS.
      apply (wp_pure_step_at m'
        (EApp (EFix body) v)
        (EApp (subst (EFix body) body) v) S).
      * apply not_value_app_fix.
      * apply step_fix_progress; exact Hvalue.
      * intros e' Hs. apply step_fix_det; auto.
      * apply (Hbody (EFix body) (VFix body) m' Hlater v Hvalue HS).
  - apply iLob.
Qed.
\end{lstlisting}
\noindent
No induction on $n$ appears in this proof; all recursive content is
hidden inside \texttt{iLob}, written once in Section~\ref{sec:prop}.
The soundness predicate $\mathtt{recursive\_spec}$ is recursive in
the function being verified, so the natural proof must invoke some
recursive reasoning principle.  A meta-level induction would have to
manufacture that recursion from outside; the internal Löb rule has
it built in.  Section~\ref{sec:imp-vs-lam} makes this structural
observation precise.

\paragraph{Distinct pre- and postconditions.}
The single-predicate \texttt{recursive\_spec} keeps the diagnostic of
Section~\ref{sec:imp-vs-lam} clean.  The mechanization also includes
the generalised form for actual program verification:
\begin{lstlisting}
Definition recursive_spec_pp (e : expr) (P Q : expr -> Prop) : iProp :=
  MkProp (fun n => forall v, value v -> P v -> wp_at n (EApp e v) Q)
         (recursive_spec_pp_mono e P Q).

Theorem wp_fix_pp body P Q :
  (forall e_self, value e_self ->
     recursive_spec_pp e_self P Q
       ⊢ recursive_spec_pp (subst e_self body) P Q) ->
  ⊤ ⊢ recursive_spec_pp (EFix body) P Q.
\end{lstlisting}
\noindent
The proof of \texttt{wp\_fix\_pp} is structurally identical to
\texttt{wp\_fix}: the same factoring through $\triangleright \varphi
\to \varphi$, the same pure-step for fix-unfolding, the same
\texttt{apply iLob}.  The locus argument of
Section~\ref{sec:imp-vs-lam} carries over without change.

\paragraph{Worked examples.}
Two short verifications exercise the framework.  The first is a
fix-based identity, observably equal to $\lambda x.x$ but forcing a
fix-unfolding at every application:
\begin{lstlisting}
Definition recursive_id : expr := EFix (ELam (EVar 0)).

Theorem recursive_id_spec (S : expr -> Prop) :
  ⊤ ⊢ recursive_spec recursive_id S.
\end{lstlisting}
\noindent
The body $\mathtt{ELam}\, (\mathtt{EVar}\, 0)$ does not mention the
self-reference, so substitution is the identity and the Lipschitz
premise becomes a non-recursive statement discharged by
\texttt{wp\_pure\_step\_at}.  The second example is the canonical
divergent fix-point:
\begin{lstlisting}
Definition bottom : expr := EFix (EVar 0).

Theorem bottom_spec (S : expr -> Prop) :
  ⊤ ⊢ recursive_spec bottom S.
Proof. apply wp_fix. intros e_self Hself. apply ientails_refl. Qed.
\end{lstlisting}
\noindent
For $\mathtt{bottom}$ the body is the bare self-reference, so
substitution gives $e_{\mathrm{self}}$ unchanged and the Lipschitz
premise collapses to reflexivity.  The conclusion --- that
$\mathtt{bottom}$ satisfies every recursive spec --- is the familiar
observation that a divergent program satisfies any partial-correctness
specification vacuously; here it appears as a one-liner because the
forcing framework absorbs the divergence into the wp's index hierarchy.
